\documentclass[11pt,a4paper]{article}
\usepackage[margin=1in]{geometry}
\usepackage{hyperref}
\usepackage{enumitem}
\usepackage{graphicx}
\usepackage{xcolor}
\usepackage{array}

% -----------------------------------------------------
% bvraghav-tiet-ucs503p-article.sty
% -----------------------------------------------------

% Variable: Lab Instructor
\makeatletter
\newcommand{\BvrSetLabInstructor}[1]{%
  \gdef\Bvr@LabInstructor{#1}}
\newcommand{\Bvr@LabInstructor}{}
\newcommand{\BvrLabInstructorValue}{\Bvr@LabInstructor}
\makeatother

% Add subtitle
\usepackage{titling}
\pretitle{\begin{center}\bfseries\fontsize{18bp}{18bp}\selectfont}
\makeatletter
\newcommand{\BvrSubtitle}[1]{%
  \posttitle{%
    \par\end{center}
    \begin{center}\large#1\end{center}
    \vskip 0.5em}%
}
\makeatother

% Hints in section titles
\newcommand{\BvrHint}[1]{%
  {\normalfont\normalsize\itshape\hfill #1}}

% -----------------------------------------------------

\title{PinPoint: A Community Driven Map for Local Discoveries}

\BvrSetLabInstructor{}

\BvrSubtitle{%
  Project Proposal for UCS503 \\
  Software Engineering \\
  Thapar Institute of Engineering and Technology
}

\author{
  Tanmay Singh, COE%
  \thanks{Roll No: \texttt{1024030338},
  Email: \texttt{tsingh2\_be24@thapar.edu}}
  \and
  Sanehbir Kaur, COE%
  \thanks{Roll No: \texttt{1024030322},
  Email: \texttt{skaur5\_be24@thapar.edu}}
}

\date{\today}

\begin{document}
\maketitle

\section{Higher Order Goal%
  \BvrHint{\ldots primary motivation}}

The project aims to provide a structured way to preserve and discover local
geographic knowledge that is often informal, temporary, personal, or difficult
to represent using conventional digital maps. The immediate engineering
objective is to build a deployable web platform where users can capture,
organize, discover, verify, and maintain such locations using geo tagged
photographs.

The project is designed as a four month software engineering effort with
emphasis on secure architecture, geospatial data management, community
workflows, maintainability, testing, and deployment.

\section{Time to Value%
  \BvrHint{\ldots development approach}}

\begin{itemize}[leftmargin=0.7cm]
    \item The project will be developed incrementally, beginning with
    authentication, maps, geo tagged locations, photographs, and
    public and private visibility.
    
    \item Moderate features such as search, filtering, community verification,
    conflict handling, and freshness tracking will follow the basic system.
    
    \item Advanced AI assisted classification, similarity detection,
    moderation, and production monitoring will be implemented after the core
    workflow is stable.
    
    \item Git based version control, automated testing, and CI/CD will be
    integrated throughout development rather than added at the end.
\end{itemize}

\section{Problem Statement}

Conventional digital maps are effective at representing established
businesses, roads, addresses, and recognized points of interest. However,
many useful locations are informal, temporary, locally known, or difficult to
represent as conventional places.

Examples include:

\begin{itemize}[leftmargin=0.7cm]
    \item Underrated food stalls and local vendors
    \item Scenic viewpoints and sunset locations
    \item Trail entrances and local shortcuts
    \item Temporary markets
    \item Public facilities
    \item Community gathering spots
    \item Lesser known local attractions
\end{itemize}

Information about these locations is commonly shared through personal
recommendations, social media, or messaging groups. Such information is
therefore difficult to discover, organize, verify, and preserve.

The problem is that useful local geographic knowledge exists, but there is no
structured community driven system specifically designed for capturing and
discovering these locations while also allowing users to maintain private
personal maps.

\section{Proposed Solution%
  \BvrHint{\ldots Software Proposition}}

\subsection{Overview}

Build a \textbf{web based} crowdsourced mapping platform called
\textbf{PinPoint} with the following components:

\begin{itemize}[leftmargin=0.7cm]
    \item \textbf{Public Maps:} users can publish useful local locations for
    discovery by the wider community.
    
    \item \textbf{My Maps:} users can maintain private maps containing
    personal locations, memories, or places that should not be publicly
    visible.
    
    \item \textbf{Geo tagged submissions:} users can upload photographs with
    geographic coordinates and descriptions.
    
    \item \textbf{AI assisted public classification:} public submissions are
    assigned tags from the approved AI generated tag suggestions.
    
    \item \textbf{Community verification:} users can add photographs, confirm
    whether locations are still available, and report incorrect information.
    
    \item \textbf{Conflict handling:} nearby or semantically similar
    locations will be identified to reduce duplicate entries.
    
    \item \textbf{Admin and Moderator dashboard:} administrators can review
    reports, flagged submissions, conflicts, and moderation activity.
\end{itemize}

\subsection{Core Workflow}

\begin{enumerate}[leftmargin=0.7cm]
    \item User signs in to the platform.
    
    \item User chooses to browse Public Maps or manage My Maps.
    
    \item User creates a location by uploading a photograph, capturing its
    geographic coordinates, and entering a description.
    
    \item User chooses whether the location belongs to a Public Map or a
    private My Map.
    
    \item For Public Maps, AI recommends tags from the approved public
    taxonomy. The user can select only from these suggested tags.
    
    \item For My Maps, the user can define custom tags and descriptions.
    
    \item The system checks nearby and similar existing public locations
    before publishing.
    
    \item Public locations can subsequently receive additional photographs,
    verification updates, reports, and freshness information.
\end{enumerate}

\subsection{Public and Private Maps}

\begin{itemize}[leftmargin=0.7cm]
    \item \textbf{Public Maps:} visible to the PinPoint community and subject
    to AI assisted tagging, moderation, duplicate detection, and community
    reporting.
    
    \item \textbf{My Maps:} private by default and accessible only to the
    owner or users explicitly granted access.
    
    \item Private maps can be shared with selected users through explicit
    access permissions.
    
    \item Private locations will not appear in public map searches or public
    APIs.
\end{itemize}

\subsection{Conflict and Duplicate Handling}

When a user attempts to add a public location near an existing location, the
system will identify possible conflicts rather than automatically creating
a duplicate.

\begin{itemize}[leftmargin=0.7cm]
    \item Display nearby or semantically similar existing locations.
    \item Allow the user to identify whether the location is the same.
    \item If it is the same, the user can add another photograph or additional
    information to the existing location.
    \item If it is different, the user can continue creating a separate
    location.
\end{itemize}

\section{Solution Approach%
  \BvrHint{\ldots Engineering focus}}

This is an engineering focused software project. The primary emphasis will
be on secure application architecture, geospatial data management, community
workflows, maintainability, and deployability rather than developing novel
AI algorithms.

\subsection{System Architecture}

The system will use a layered web architecture:

\begin{itemize}[leftmargin=0.7cm]
    \item \textbf{Frontend:} React with TypeScript and Tailwind CSS.
    \item \textbf{Backend:} FastAPI with Python.
    \item \textbf{Database:} PostgreSQL with PostGIS and pgvector.
    \item \textbf{Object Storage:} AWS S3 for photographs.
    \item \textbf{Caching and Background Processing:} Redis and Celery.
    \item \textbf{Reverse Proxy:} Nginx.
    \item \textbf{Deployment:} Vercel for the frontend and AWS EC2 with
    Docker Compose for backend services.
\end{itemize}

\subsection{Geospatial Management}

Leaflet and OpenStreetMap will provide interactive map functionality.
Nominatim will be used for geocoding and reverse geocoding.

PostGIS will store geographic coordinates and support proximity queries used
for location discovery and duplicate and conflict detection.

\subsection{AI Assisted Public Maps}

AI will be used as an assistance and moderation layer.

\begin{itemize}[leftmargin=0.7cm]
    \item \textbf{Gemini API:} recommends public categories and tags from
    user provided descriptions.
    
    \item \textbf{Sentence Transformers:} generate semantic embeddings for
    location descriptions and related information.
    
    \item \textbf{pgvector:} performs similarity searches for potentially
    duplicate or conflicting locations.
    
    \item AI assisted moderation will identify potentially inappropriate,
    spam, or suspicious public submissions.
    
    \item AI will assist moderation but will not override application rules
    or administrator decisions.
\end{itemize}

\subsection{Authentication and Security}

Authentication will be mandatory for all users.

\begin{itemize}[leftmargin=0.7cm]
    \item JWT and OAuth2 for authentication.
    \item Argon2id for password hashing.
    \item Ownership and authorization checks for private maps.
    \item Explicit sharing permissions for selected users.
    \item Input validation through Pydantic.
    \item API rate limiting.
    \item Secure access to stored photographs through controlled URLs.
    \item HTTPS and Nginx for secure communication.
\end{itemize}

\subsection{Testing and Validation}

The system will be validated through multiple levels of testing:

\begin{itemize}[leftmargin=0.7cm]
    \item \textbf{Unit Testing:} individual backend functions and services
    using Pytest.
    
    \item \textbf{Integration Testing:} API, database, authentication,
    storage, and background service interactions.
    
    \item \textbf{End to End Testing:} complete user workflows using
    Playwright.
    
    \item \textbf{Security Testing:} authentication, authorization, input
    validation, rate limiting, and dependency and container security checks.
    
    \item \textbf{User Acceptance Testing:} testing public and private map
    workflows, location creation, discovery, verification, and reporting
    with representative users.
\end{itemize}

\subsection{CI/CD and Development Workflow}

The project will use Git and GitHub for collaborative development.

\begin{itemize}[leftmargin=0.7cm]
    \item Feature branches will be used for individual features.
    \item Pull requests will be used for code review and merging.
    \item GitHub Actions will automatically run linting, type checking,
    tests, security checks, and Docker builds.
    \item Docker and Docker Compose will provide reproducible development and
    deployment environments.
    \item Releases will use version tags and semantic versioning.
\end{itemize}

\section{Evaluation Criterion%
  \BvrHint{\ldots measurable and attributable}}

The system will be evaluated using metrics connected to the primary workflows.

\subsection{Primary Evaluation Metrics}

\begin{itemize}[leftmargin=0.7cm]
    \item \textbf{Location creation success rate:} percentage of valid
    submissions successfully stored and displayed.
    
    \item \textbf{Public and private authorization correctness:} percentage
    of tested private locations that remain inaccessible through unauthorized
    users and public APIs.
    
    \item \textbf{Duplicate detection effectiveness:} percentage of known
    duplicate or nearby conflicting locations correctly surfaced for review.
\end{itemize}

\subsection{Secondary Evaluation Metrics}

\begin{itemize}[leftmargin=0.7cm]
    \item API response time for common map and search operations.
    \item Successful image upload and retrieval rate.
    \item Community verification and reporting response rate.
    \item Freshness of public location information.
    \item Application availability during the deployment period.
    \item End to end workflow success rate.
\end{itemize}

\section{Scalability%
  \BvrHint{\ldots with foresight}}

The architecture will support growth in users, locations, photographs, and
concurrent requests.

\begin{enumerate}[leftmargin=0.7cm]
    \item \textbf{Application scalability:}
    \begin{itemize}[leftmargin=0.7cm]
        \item Separate frontend and backend services.
        \item Stateless FastAPI APIs.
        \item Redis based caching.
        \item Celery based asynchronous processing.
        \item Pagination and indexed database queries.
        \item Background execution of expensive AI operations.
    \end{itemize}

    \item \textbf{Infrastructure scalability:}
    \begin{itemize}[leftmargin=0.7cm]
        \item Managed PostgreSQL through AWS RDS.
        \item AWS S3 for scalable photograph storage.
        \item CloudFront for image delivery and caching.
        \item Docker based deployment.
        \item CI/CD for repeatable releases.
    \end{itemize}
\end{enumerate}

\section{Engine Availability Heuristic}

The project uses mature and widely adopted technologies:

\begin{itemize}[leftmargin=0.7cm]
    \item React and TypeScript with Tailwind CSS for frontend development.
    \item FastAPI and Python for backend APIs.
    \item PostgreSQL with PostGIS and pgvector for relational, geospatial,
    and vector data.
    \item Leaflet and OpenStreetMap for mapping.
    \item Nominatim for geocoding.
    \item AWS S3 and CloudFront for image storage and delivery.
    \item Redis and Celery for caching and background jobs.
    \item Gemini API and Sentence Transformers for AI functionality.
    \item Docker and Docker Compose for containerization.
    \item Pytest and Playwright for automated testing.
    \item GitHub Actions for CI/CD.
    \item Sentry and AWS CloudWatch for monitoring.
\end{itemize}

These technologies allow the team to focus on the application's core
workflow, security, geospatial functionality, and community features rather
than building infrastructure from scratch.

\section{Project Scope and Deliverables%
  \BvrHint{\ldots four month implementation}}

\subsection{Initial Deliverable: Basic}

\begin{itemize}[leftmargin=0.7cm]
    \item Project setup and Git based development workflow.
    \item Mandatory user registration and authentication.
    \item React based map interface.
    \item Leaflet and OpenStreetMap integration.
    \item Creation and viewing of geo tagged locations.
    \item Photograph upload to AWS S3.
    \item Basic public and private location visibility.
    \item Creation and management of My Maps.
    \item PostgreSQL and PostGIS integration.
\end{itemize}

\subsection{Subsequent Deliverable: Moderate}

\begin{itemize}[leftmargin=0.7cm]
    \item Public categories and filtering.
    \item Location search and detail cards.
    \item Multiple photographs for locations.
    \item Nearby location and duplicate suggestions.
    \item Community verification and reporting.
    \item Location freshness and status tracking.
    \item Private map sharing with selected users.
    \item Admin and Moderator dashboard.
    \item Redis and Celery background processing.
\end{itemize}

\subsection{Advanced Deliverable}

\begin{itemize}[leftmargin=0.7cm]
    \item Gemini based public tag recommendation.
    \item Sentence Transformer embeddings with pgvector.
    \item AI assisted duplicate and similarity detection.
    \item AI assisted spam and content moderation.
    \item Advanced conflict and freshness handling.
    \item Security hardening and automated security checks.
    \item Production CI/CD and monitoring.
\end{itemize}

\subsection{Four Month Timeline}

\begin{center}
\renewcommand{\arraystretch}{1.25}
\begin{tabular}{|p{2.2cm}|p{5.2cm}|p{5.2cm}|}
\hline
\textbf{Period} & \textbf{Development Phase} & \textbf{Deliverable} \\
\hline
Weeks 1 to 4 &
Requirements, architecture, database design, project setup, authentication &
SRS, architecture, repository, basic authentication \\
\hline
Weeks 5 to 8 &
Maps, geolocation, image upload, public and private maps &
Basic working PinPoint MVP \\
\hline
Weeks 9 to 12 &
Search, filters, conflicts, community verification, freshness, admin dashboard &
Moderate feature complete system \\
\hline
Weeks 13 to 15 &
AI tagging, embeddings, duplicate detection, moderation &
Advanced AI assisted functionality \\
\hline
Week 16 &
Testing, security checks, deployment, monitoring, documentation &
Final release, testing report, presentation \\
\hline
\end{tabular}
\end{center}

\section{Risks and Mitigations}

\begin{itemize}[leftmargin=0.7cm]
    \item \textbf{Privacy risk:} enforce strict ownership and access control
    rules for private maps and photographs.
    
    \item \textbf{Incorrect AI classification:} public tags will be selected
    only from approved AI suggested tags and remain subject to moderation.
    
    \item \textbf{Duplicate locations:} combine geographic proximity and
    semantic similarity to identify possible duplicates for user review.
    
    \item \textbf{Outdated information:} maintain verification timestamps,
    status updates, and community reporting.
    
    \item \textbf{Spam and inappropriate content:} combine AI assisted
    moderation, rate limiting, reporting, and administrator review.
    
    \item \textbf{Project scope:} implement features in Basic, Moderate, and
    Advanced stages so that the core system is completed before advanced
    functionality.
    
    \item \textbf{Deployment issues:} use Docker, automated tests, CI/CD,
    Sentry, and CloudWatch to make deployment repeatable and observable.
\end{itemize}

\section{Summary}

PinPoint proposes a deployable web platform for capturing and discovering
local geographic knowledge that is difficult to represent through
conventional maps. The system combines geo tagged photographs, public
community maps, private personal maps, geospatial search, community
verification, freshness tracking, and AI assisted public classification and
moderation.

The project emphasizes software engineering principles through modular
architecture, secure authentication and authorization, automated testing,
containerization, CI/CD, monitoring, and managed cloud infrastructure. The
four month implementation will progress from a basic mapping system to
moderate community features and finally advanced AI assisted functionality.

\end{document}